\documentclass[10pt]{article}

\usepackage[T1]{fontenc}
\usepackage[utf8]{inputenc}
\usepackage{lmodern}
\usepackage{microtype}
\usepackage[a4paper,margin=17mm,headheight=14pt]{geometry}
\usepackage{booktabs}
\usepackage{tabularx}
\usepackage{array}
\usepackage{float}
\usepackage{enumitem}
\usepackage{xcolor}
\usepackage{hyperref}
\usepackage{fancyhdr}
\usepackage{titlesec}
\usepackage{siunitx}

\definecolor{navy}{HTML}{15324A}
\definecolor{blue}{HTML}{2F6F9F}
\definecolor{pale}{HTML}{EAF2F8}
\definecolor{green}{HTML}{DDEFE4}
\definecolor{amber}{HTML}{FFF1CC}
\definecolor{grey}{HTML}{5D6973}

\hypersetup{
  colorlinks=true,
  linkcolor=navy,
  urlcolor=blue,
  pdftitle={Public-Data Hedging Benchmark Summary},
  pdfauthor={Quantitative Research}
}

\pagestyle{fancy}
\fancyhf{}
\lhead{Public-Data Hedging Benchmark}
\rhead{13 October 2026}
\cfoot{\thepage}
\renewcommand{\headrulewidth}{0.3pt}

\titleformat{\section}{\large\bfseries\color{navy}}{\thesection}{0.6em}{}
\titleformat{\subsection}{\normalsize\bfseries\color{blue}}{\thesubsection}{0.5em}{}
\titlespacing*{\section}{0pt}{1.1em}{0.45em}
\titlespacing*{\subsection}{0pt}{0.8em}{0.3em}

\setlist[itemize]{leftmargin=1.4em,itemsep=0.15em,topsep=0.25em}
\setlength{\parindent}{0pt}
\setlength{\parskip}{0.35em}
\renewcommand{\arraystretch}{1.16}
\raggedbottom

\newcommand{\statuscomplete}{\colorbox{green}{\strut\textbf{complete}}}
\newcommand{\statusready}{\colorbox{pale}{\strut\textbf{ready}}}
\newcommand{\statuspending}{\colorbox{amber}{\strut\textbf{pending}}}
\newcommand{\model}[1]{\textbf{#1}}

\begin{document}

\begin{center}
  {\LARGE\bfseries\color{navy} Public-Data Hedging Benchmark}\par
  \vspace{0.2em}
  {\large Approach, current results, and model catalogue}\par
  \vspace{0.35em}
  {\color{grey}Status date: 13 October 2026}
\end{center}

\vspace{0.4em}
\noindent\colorbox{pale}{%
  \parbox{\dimexpr\textwidth-2\fboxsep\relax}{%
  \textbf{Executive status.}
  The model-neutral daily benchmark is operational on a normalized,
  original-data-free public-candidate panel. Two methods have complete
  out-of-sample results over all 611 eligible 2023--2025 origins and all 20
  daily digital targets. SANOS--LV reduces aggregate pure-hedge RMSE by
  \textbf{56.5\%} versus the unhedged claims; historical statistical SANOS
  reduces it by \textbf{17.4\%}. Robust covariance, GROMM, VolGAN, and
  AD--Seq--Vol have passed their current implementation gates but do not yet
  have comparable full-production results. Weekly mechanics and pricing pass;
  weekly production evaluation remains outstanding.}}

\section{Approach}

\subsection{Data boundary and common market}

The sole runtime input is the distributable-candidate dataset
\texttt{spx-2022-2025-hedging-public-candidate-v1}. It contains 891 retained
dates from 3 January 2022 to 29 August 2025 and 2,014,740 contemporaneously
eligible quote-side option coordinates. It contains no original WRDS price,
midpoint, identifier, volume, open interest, vendor Greek, or vendor implied
volatility. Instead it releases normalized fitted SANOS and DLV states,
separate reconstructable SANOS and DLV model prices, Black implied
volatilities, forward/spot ratios, discount factors, normalized half-spreads,
and forward-adjusted spot returns.

The common executable market mark is the released SANOS price. A purchase pays
SANOS plus one half-spread and a sale receives SANOS minus one half-spread.
Spot costs one basis point of traded notional per transaction; cash settlement
is free. The release is technically classified as \texttt{public\_candidate};
redistribution still requires a separate rights decision.

\subsection{Out-of-sample protocol}

\begin{itemize}
  \item Model initialization, architecture validation, and target sizing use
  2022 only. Every rolling estimator uses the previous 252 exact one-session
  transitions and no future observation.
  \item Daily evaluation runs from 17 January 2023 through 28 August 2025:
  611 eligible origins, 20 targets per origin, and 12,220 claim/origin records.
  \item The targets are cash digitals at 1D and 3D: ATM call and put plus OTM
  NMS levels $\pm0.25,\pm0.50,\pm0.75,\pm1.00$. Quantities are fixed from
  2022 so each target has one maturity-matched ATM spot-risk unit.
  \item The hedge universe is normalized spot plus one OTM option per eligible
  strike inside $(-2,2)$ NMS, using expiries up to five business days (or the
  first later expiry when five days is absent).
  \item Weekly evaluation starts an overlapping five-session episode on every
  eligible date. Frozen weekly fits once and keeps its inception contract book;
  restarted weekly refits daily. Both may react to the realized state.
\end{itemize}

\subsection{Common accounting and objectives}

Every method starts from the same market claim mark and is evaluated on the
same realized next-date SANOS mark or cash payoff. The report separates
\emph{pricing loss}, \emph{pure hedge loss}, and their sum, \emph{economic
loss}. Buyer and seller indifference prices are additional model outputs, not
substitutes for this common-price comparison.

Supported risk criteria are mean--variance, mean--volatility, entropy OCE, and
power-exponential OCE. Learned or statistical policies use the hedge of the
claim relative to the corresponding zero-claim policy; generated models also
fit a direct pure policy so transaction costs are charged on the executable
net position rather than inferred from differences of absolute-policy costs.

\section{Current Results}

\subsection{Completed daily production runs}

Table~\ref{tab:risk} uses the pure-hedge loss in inception-spot units. Aggregate
RMSE is the square root of the equally weighted mean claim-level MSE; every
claim contributes 611 observations. ``Reduction'' is relative to the matching
unhedged RMSE. Expected shortfall is the equally weighted claim-level 95\%
expected shortfall.

\begin{table}[H]
\centering
\small
\caption{Out-of-sample daily risk across the completed 611-origin runs.}
\label{tab:risk}
\begin{tabular}{llrrrrr}
\toprule
Model & Targets & Unhedged RMSE & Pure RMSE & Reduction & Economic RMSE & ES$_{95}$ \\
\midrule
SANOS--LV & all 20 & 0.01632 & 0.00710 & 56.5\% & 0.00710 & 0.01351 \\
SANOS--LV & 1D & 0.01244 & 0.00862 & 30.7\% & 0.00862 & 0.01329 \\
SANOS--LV & 3D & 0.01943 & 0.00514 & 73.6\% & 0.00514 & 0.01373 \\
\addlinespace
Historical SANOS & all 20 & 0.01632 & 0.01347 & 17.4\% & 0.01347 & 0.02571 \\
Historical SANOS & 1D & 0.01244 & 0.01116 & 10.3\% & 0.01116 & 0.02234 \\
Historical SANOS & 3D & 0.01943 & 0.01544 & 20.6\% & 0.01544 & 0.02908 \\
\bottomrule
\end{tabular}
\end{table}

The SANOS--LV hedge is materially stronger on this first completed comparison,
especially for the remaining-life 3D claims. Its aggregate pure RMSE is 47.3\%
lower than historical SANOS. Pricing loss is negligible at aggregate scale:
SANOS--LV economic RMSE is 0.007098 versus pure RMSE 0.007096; historical SANOS
uses the common price and therefore has zero pricing contribution by design.

\begin{table}[H]
\centering
\small
\caption{Average trading diagnostics across all 20 targets and 611 origins.}
\label{tab:trading}
\begin{tabular}{lrrrrrr}
\toprule
Model & Mean cost & Entry turnover & Nonzero actions & Cap hits & Nonconverged & Records \\
\midrule
SANOS--LV & 0.001011 & 2.670 & 5.0 & 0 & 0 & 12,220 \\
Historical SANOS & 0.001019 & 0.0738 & 181.6 & 0 & 0 & 12,220 \\
\bottomrule
\end{tabular}
\end{table}

The similar average cost hides different portfolio geometry. SANOS--LV uses a
sparse exact eight-factor hedge plus spot; historical SANOS spreads small
positions across much of the 200-coordinate option book. Turnover and action
count therefore should not be read interchangeably.

\subsection{Evidence level across models}

\begin{table}[H]
\centering
\small
\caption{Current evidence. ``Readiness'' is not a performance estimate.}
\begin{tabularx}{\textwidth}{>{\raggedright\arraybackslash}p{28mm}p{19mm}p{25mm}X}
\toprule
Model & Measure & Daily status & Evidence \\
\midrule
SANOS--LV & Q price / realized path & \statuscomplete & 611 origins, 12,220 records, zero cap hits and zero nonconvergence. \\
Historical SANOS & P & \statuscomplete & 611 origins, 12,220 records, zero cap hits and zero nonconvergence. \\
Robust covariance & P & \statusready & Full 20-target one-origin mean--volatility gate passed for $\kappa\in\{0.1,1,2,3\}$; production pending. \\
GROMM & P and Q & \statusready & Rolling artifacts, daily/weekly scenarios, four objectives, direct Q-shadow pure hedge, and quant review pass; managed-GPU fit and production pending. \\
VolGAN & P & \statusready & Public-architecture implementation, restartable rolling updates, daily/weekly adapters, and smoke tests pass; production pending. \\
AD--Seq--Vol / FT & P & \statusready & Full-geometry managed-GPU origin: 4,096 scenarios, 20 targets, 158.3 seconds, 2,636 MiB VRAM, no cap hit or nonconvergence; production pending. \\
\bottomrule
\end{tabularx}
\end{table}

\subsection{Supporting calibration and weekly evidence}

The rolling rank-16 DLV innovation fit explains 86.50\% of variance at the
first production origin. The compatible physical drift passes its precommitted
holdout test: 1.87\% validation improvement, paired gain statistic 6.84, and
daily risk-price norm 0.371. These are model-calibration diagnostics, not hedge
performance results.

The weekly engine has 528 eligible origins through 22 August 2025 and 13 common
targets: ten five-day digitals, realized variance, a 95\% end-of-day down-and-out
ATM put, and a 105\% end-of-day up-and-out ATM call. The common DLV transition
test has maximum mass and martingale errors below $3.6\times10^{-14}$ across an
eight-date check. Daily, frozen-weekly, and restarted-weekly adapters exist for
all model families, but no complete weekly comparison is yet reported.

\section{Model Catalogue}

\subsection{SANOS--LV reference hedge}

SANOS supplies the common market interpolation and digital price. DLV local
volatility supplies the model price and Greeks. From each trailing 252-day
window, the method estimates one front-spot-projected DLV factor and seven
residual PCA factors. Central DLV-squared bumps give factor vegas; a linear
program selects the minimum predictable-cost option portfolio that exactly
neutralizes them, and spot removes residual delta. One-day claims are
delta-only; three-day claims use at most eight option hedges. The method is
deterministic, sparse, and the strongest completed baseline.

\subsection{Independent robust covariance}

This baseline deliberately avoids a SANOS scenario generator. It replays the
current relative option book over 252 historical public transitions, fits
robust scales and a regularized Tyler scatter matrix, and estimates each target
cross-moment separately by robust regression. Expected returns are HAC-screened
in covariance eigenportfolios. Mean--variance and mean--volatility hedges use
spread-weighted lasso costs; the latter is checked analytically for recession
directions. The implementation gate passes, but only one production origin
exists and it must not be compared with Table~\ref{tab:risk}.

\subsection{Historical statistical SANOS}

Historical SANOS reconstructs 252 same-coordinate option and target returns
from released SANOS states and uses those empirical rows directly as scenarios.
It supports all four risk criteria. Long, short, and zero policies determine
indifference prices; a separate drift-neutral policy gives the executable pure
hedge. Mean--variance uses exact cost-lasso coordinate descent and nonlinear
objectives use analytic-gradient L-BFGS-B. The model is inexpensive and fully
rolling, but its completed mean--volatility $\kappa=1$ result trails SANOS--LV.

\subsection{GROMM}

GROMM is the native generative risk-neutral option-market model: a 200-strike
JAX market fit, 16-factor AR(1) log-DLV-squared dynamics, independent 5\% Black
background, and eight Brownian substeps per business day. Dated artifacts are
bound to exactly 252 public transitions and must pass compatible-P drift gates.
It generates Q and accepted compatible-P scenarios. An affine-plus-deep LSTM
supports all four objectives; the executable pure policy uses conditional
Q-shadow gains at each realized P state, removing predictable physical drift
without fitted holdout centering. CPU integration and quant review pass;
rolling managed-GPU production remains outstanding.

\subsection{VolGAN}

VolGAN is a historical-P surface generator on a fixed $9\times11$ Black-IV
grid reconstructed from public SANOS states. It conditions on two lagged
returns, trailing 21-day realized volatility, and the current surface. The
published two-layer Softplus architecture uses 32 noise dimensions, a 16-unit
hidden layer, and strike/maturity gradient penalties. Weekly paths recursively
update surface and return state. The implementation is project-native and uses
no proprietary surfaces. Smoke, restart, and rolling warm-start gates pass;
full production does not yet exist.

\subsection{AD--Seq--Vol and arbitrage-penalized fine tuning}

AD--Seq--Vol is a causal sequence diffusion model on 22-day histories of the
same $9\times11$ surface plus forward-adjusted returns. Its eight-layer,
width-256, four-head transformer predicts the clean sequence through a 200-step
cosine DDPM. AD--Seq--Vol--FT freezes the EMA base and trains a zero-initialized
rank-four adapter with calendar, call-spread, and butterfly penalties. Both are
historical-P generators and supply indifference prices rather than Q
expectations. Exact restart, rolling warm-start, daily/weekly adapters, and one
full-geometry GPU origin pass; production and batch-invariance remain pending.

\section{Interpretation and Next Evidence}

The benchmark now has a reproducible common market, accounting system, target
catalogue, cost model, adapters, and two complete daily baselines; it does not
yet establish a final ranking. Next come complete covariance, GROMM P/Q, and
generated-model daily runs; production batch/restart checks; and the frozen-
versus-restarted weekly comparison with overlap-aware inference.

\vspace{0.4em}
\hrule
\vspace{0.35em}
{\footnotesize
\textbf{Provenance.}
The completed results use public dataset manifest SHA-256
\texttt{58741552e0f6...12cf903007a}. SANOS--LV run identity begins
\texttt{2af7a4f64df4}; historical-SANOS run identity begins
\texttt{64adf168a213}. All headline numbers are derived from immutable
\texttt{hedging-daily-results-v2} artifacts. Readiness statements come from
the dated reports and independent reviews alongside the implementation.}

\end{document}
